\documentclass[10pt,a4paper]{amsart}
\usepackage[margin=19mm]{geometry}
\usepackage{amsmath,amssymb,mathtools}
\usepackage[scr=rsfs,cal=euler]{mathalfa}
\usepackage{xfrac}
\usepackage[unicode,hidelinks,bookmarks=false]{hyperref}
\newcommand{\ie}{i.e.\ }
\newcommand{\cf}{cf.\ }
\newcommand{\resp}{respectively\ }
\newcommand{\Subsection}{\subsection}
\providecommand{\nobreakdash}{\nobreak\hbox{-}}
\setlength{\emergencystretch}{2em}
\multlinegap0pt
\usepackage{amssymb}
\usepackage{enumerate}
\usepackage{dsfont}
\usepackage{bm}
\usepackage{xcolor}
\newtheorem{lemma}{Lemma}
\def \N {\mathbb{N}}
\def \R {\mathbb{R}}
\def \Z {\mathbb{Z}}
\def\d{{\,{\rm d}}}
\renewcommand{\leq}{\leqslant}
\title{Counting solutions to the quadratic determinant equation}
\author{Jonathan Chapman, Akshat Mudgal}
\date{Source: arXiv:2605.15434v1 (CC BY 4.0)}
\begin{document}
\maketitle

\noindent\emph{Source and licence.} Counting solutions to the quadratic determinant equation. Version arXiv:2605.15434v1 (CC BY 4.0), distributed by arXiv under the Creative Commons Attribution 4.0 International licence, http://creativecommons.org/licenses/by/4.0/, as recorded in the arXiv licence metadata for this exact version and shown on the abstract page https://arxiv.org/abs/2605.15434v1. Full licence notice: \texttt{licenses/arxiv-2605.15434-CC-BY-4.0.txt}.\par\smallskip

\noindent\textit{Editorial scope.} Counted unit: the article's lemma lem3.2 (source0.tex lines 516--519). The article's complete original proof of it is delivered with it (source0.tex lines 520--571, one \texttt{proof} environment of 52 lines). No mathematics is written, repaired or re-derived: the statement and the proof are the article's own lines, in the article's own notation, and no retained line is substituted or re-flowed.  Retained uncounted as the source-local context the proof needs: lines 417--419; lines 453--456; lines 511--513.  Nothing was omitted from any retained span, so the package carries no omission record.  No retained line contains a citation, so the wrapper carries no bibliography and no bibliography entry.  No retained line contains a figure environment, an image-inclusion command or drawing code, so no image is delivered and the package declares no asset dependency.  Editorial formatting only, in addition: an AMS \texttt{amsart} preamble that restates the article's own declarations and macros used by the retained lines, namely the environments \texttt{lemma} on separate counters, together with the standard packages those lines need.  The retained environments print in this document as Lemma 1 in order of appearance, whereas the article prints the same items with the numbering of its own full document.  The complete native source of the article as distributed in the arXiv e-print for this exact version is bundled unchanged as \texttt{sources/200652/source0.tex}.
\begin{equation}\label{eqn3.5}
   0\leqslant  F_N(\alpha,\beta,h/d) \leqslant \frac{N}{\max\{\alpha,\beta\}}.
\end{equation}

\begin{lemma} \label{lem3.1}
    Let $N\in\N\setminus\{1\}$ and $h\in\Z$ with $0\leqslant h\leqslant 2N^2$. Let $d,k\in[N]$ with $d\mid h$ and $k\leqslant N/d$. We have
    \[  \sum_{1 \leq u,v \leq N/dk} \int_{u}^{u+1} \int_{v}^{v+1} |  F_N(u,v,h/dk) -   F_N(\alpha, \beta, h/dk) |  \d \alpha \d \beta   \ll N (\log N)^2.  \]
\end{lemma}

    \begin{equation}\label{eqn3.8}
        \mu( (0, \beta] \cap (\lambda - (0, \alpha]  ) ) \leqslant\min\{\alpha,\beta,\sqrt{\alpha\beta}\} \qquad(0\leqslant\alpha,\beta\leqslant 1),
    \end{equation}

\begin{lemma} \label{lem3.2}
    Let $N,h,d,k$ be as in Lemma \ref{lem3.1}. Writing $h = \lambda N^2$, we have
    \[ S_N(h/dk) = \frac{N^2}{dk} J(\lambda)  + O(N(\log N)^2 ). \]
\end{lemma}

\begin{proof}
Recalling \eqref{eqn3.5} and applying a change of variables, we see that
\begin{align*}
S_N(h/dk) & = \frac{N}{dk} \int_{dk/N}^1 \int_{dk/N}^1 F_N( \alpha, \beta , h/N)  \d \beta \d \alpha \\
& = \frac{N}{dk} \int_{dk/N}^1 \int_{dk/N}^1 \max\left\lbrace 0,\min\left\{\frac{N}{\alpha} ,  \frac{h}{N\alpha\beta} - \frac{1}{\beta} \right\}-\max\left\{ \frac{1}{\alpha}, \frac{h}{N\alpha\beta}-\frac{N}{\beta} \right\}\right\rbrace  \d \beta \d \alpha \\
& = \frac{N^2}{dk} \int_{dk/N}^1 \int_{dk/N}^1 \frac{1}{\alpha \beta} \max \left\{  0, \min\left\{ \beta, \frac{h}{N^2}-\frac{\alpha}{N} \right\} - \max \left\{ \frac{\beta}{N} , \frac{h}{N^2}-\alpha\right\} \right\}  \d \beta \d \alpha.
\end{align*}
Now observe that for any $t,x_1,x_2,y_1,y_2\in\R$ with $x_1 \leqslant y_1$ and $x_2 \leqslant y_2$ we have
\[ \mu([x_1, y_1] \cap (t - [x_2, y_2]) ) = \max \{ 0, \min\{ y_1, t - x_2\} - \max\{ x_1, t - y_2 \} \} .\]
Recalling that $h=\lambda N^2$, we therefore infer
\[ S_N(h/dk) = \frac{N^2}{dk} \int_{dk/N}^1 \int_{dk/N}^1 \frac{1}{\alpha \beta} \ \mu \left( \left[ \frac{\beta}{N}, \beta \right] \cap \left(\lambda - \left[ \frac{\alpha}{N}, \alpha \right] \right) \right)  \d \beta \d \alpha.  \]
Furthermore, in view of the bound
\[ \frac{N^2}{dk} \int_{dk/N}^1 \int_{dk/N}^1  \frac{1}{\alpha N} \ d \beta d \alpha \ll \frac{N}{dk} ( \log N + \log (dk) ) \ll N \log N, \]
we may write
\begin{equation*}
    S_N(h/dk) = \frac{N^2}{dk} \int_{dk/N}^1 \int_{dk/N}^1 \frac{1}{\alpha \beta} \ \mu \left( \left[ 0, \beta \right] \cap \left(\lambda - \left[ 0, \alpha \right] \right) \right)  \d \beta \d \alpha +O(N\log N).
\end{equation*}
If we extend the range of both integrals from $[dk/N,1]$ to $(0,1)$, then
\begin{equation*}
    S_N(h/dk) = \frac{N^2}{dk} \int_{0}^1 \int_{0}^1 \frac{1}{\alpha \beta} \ \mu \left( \left[ 0, \beta \right] \cap \left(\lambda - \left[ 0, \alpha \right] \right) \right)  \d \beta \d \alpha +O(I_1+I_2+N\log N),
\end{equation*}
where
\begin{align*}
   I_1 &= \frac{N^2}{dk}\int_{dk/N}^1\int_0^{dk/N} \frac{1}{\alpha \beta} \ \mu( [0, \beta] \cap (\lambda - [0, \alpha] ) )\d \beta \d\alpha,\\
   I_2 &= \frac{N^2}{dk} \int_0^{dk/N} \int_{0}^1 \frac{1}{\alpha \beta} \ \mu( [0, \beta] \cap (\lambda - [0, \alpha] ) )\d \beta \d\alpha.
\end{align*}
Hence, to complete the proof, it only remains to show that $I_1, I_2=O(N(\log N)^2)$.

The crude bound \eqref{eqn3.8} is enough to establish
\begin{equation*}
    I_1\leqslant \frac{N^2}{dk}\int_{dk/N}^1 \frac{1}{\alpha}\int_0^{dk/N} 1 \d\beta\d\alpha =N\log(N/dk) \leqslant N\log N.
\end{equation*}
We now consider $I_2$. Recall that
\begin{equation*}
    \mu( [0, \beta] \cap (\lambda - [0, \alpha] ) ) = \mu([\max\{0,\lambda -\alpha\}, \min\{\beta,\lambda\}]).
\end{equation*}
Let $\ell = \min\{ \lambda,dk/N\}$. We may assume $\ell>0$, as otherwise $I_2=0$ and we are done. Thus, 
\begin{align*}
    \int_0^\ell \int_{0}^1 \frac{1}{\alpha \beta} \ \mu( [0, \beta] \cap (\lambda - [0, \alpha] ) )\d\beta \d \alpha  &= \int_0^\ell \frac{1}{\alpha}\left(\int_{\lambda -\alpha }^\lambda \frac{\alpha +\beta-\lambda}{\beta}\d\beta +\int_\lambda^1\frac{\alpha}{\beta}\d\beta \right)\d\alpha\\
    &= \ell(1-\log(\lambda)) + \lambda\int_0^{\ell/\lambda}\left(\frac{1}{t} - 1\right)\log(1-t)\d t.
\end{align*}
Observe that the function $t\mapsto (t^{-1}-1)\log(1-t)$ is continuous on $(0,1)$ and has finite limits as $t\to 0^+$ or $t\to 1^-$. Consequently, the integrand appearing in the final integral above is uniformly bounded by some absolute constant. Since we are in the case when $\lambda\neq 0$, we know that $N^{-2}\leqslant \lambda \leqslant 2$, and so
\begin{equation*}
    \int_0^\ell \int_{0}^1 \frac{1}{\alpha \beta} \ \mu( [0, \beta] \cap (\lambda - [0, \alpha] ) )\d\beta \d \alpha \ll \ell(1+|\log(\lambda)|)\ll \frac{dk}{N}\log(N).
\end{equation*}
If $\lambda\geqslant dk/N$, then $\ell=dk/N$ and this bound leads to the desired estimate $I_2=O(N\log N)$. If instead $0<\lambda<dk/N\leqslant 1$, then $\ell=\lambda\geqslant N^{-2}$ and the previous bound implies that
\begin{align*}
    I_2 + O(N\log N) &= \frac{N^2}{dk}\int_\lambda^{dk/N}\frac{1}{\alpha}\left(\int_0^{\lambda}\frac{\beta}{\beta}\d\beta+\int_\lambda^1\frac{\lambda}{\beta}\d\beta\right)\d\alpha\\
    &\ll \frac{\lambda N^2}{dk}(1+|\log(\lambda)|^2) \ll N(\log N)^2,
\end{align*}
as required. 
\end{proof}

\end{document}
